\documentclass[11pt,a4paper]{amsart}

\usepackage[T1]{fontenc}
\usepackage{amsmath,amssymb,amsthm,mathtools}
\usepackage{enumitem}
\usepackage{microtype}
\usepackage[margin=2.7cm]{geometry}
\usepackage{xcolor}
\definecolor{linkcol}{rgb}{0.10,0.20,0.50}
\usepackage[colorlinks=true,linkcolor=linkcol,citecolor=linkcol,urlcolor=linkcol]{hyperref}
\makeatletter
\g@addto@macro\UrlBreaks{\do\_\do\.}
\makeatother

\theoremstyle{plain}
\newtheorem{theorem}{Theorem}[section]
\newtheorem{proposition}[theorem]{Proposition}
\newtheorem{lemma}[theorem]{Lemma}
\newtheorem{corollary}[theorem]{Corollary}
\theoremstyle{definition}
\newtheorem{definition}[theorem]{Definition}
\newtheorem{example}[theorem]{Example}
\newtheorem{remark}[theorem]{Remark}

\newcommand{\CC}{\mathbb C}
\newcommand{\Cu}{C_u^*}
\newcommand{\Ce}{C_e^*}
\newcommand{\Esc}{\mathfrak E}
\newcommand{\Def}{\mathfrak D}
\newcommand{\Corr}{\mathbf{Corr}_{C^*}}
\newcommand{\OSq}{\mathbf{OS}_{\mathrm q}}
\newcommand{\id}{\operatorname{id}}
\newcommand{\Ideal}{\operatorname{Ideal}}
\newcommand{\coker}{\operatorname{coker}}
\newcommand{\graph}{\operatorname{graph}}
\newcommand{\ran}{\operatorname{ran}}
\emergencystretch=2em

\begin{document}

\title[The Shilov correlation law]{The Shilov correlation law for operator-system quotients}

\author{M.~Haufschild}
\thanks{Independent researcher. Email: \texttt{haufschildmichael@gmail.com}.}

\date{4 August 2026}

\subjclass[2020]{Primary 46L07, 47L25; Secondary 46M18, 18G40, 46L30}
\keywords{operator system, Shilov boundary, $C^*$-envelope, complete order quotient,
subdirect product, cohomology of diagrams, derived inverse limit}

\begin{abstract}
The $C^*$-envelope is not functorial for complete order quotients.  We prove that this failure has
a canonical exact completion.  If $q:S\twoheadrightarrow T$ is a complete order quotient,
$\alpha:\Cu(S)\twoheadrightarrow\Cu(T)$ is the induced map of universal covers, and $J_S,J_T$
are the Shilov ideals, set
\[
 M_q=J_T+\alpha(J_S),\qquad Q_q=\Cu(T)/M_q .
\]
The image of the universal graph is a closed subdirect product
\[
 \Gamma_q=\Ce(S)\times_{Q_q}\Ce(T).
\]
Thus every quotient has a canonical envelope correlation even when it has no envelope map.
The two quotient faces of the minimal join $M_q$ are precisely the escape and deficit algebras,
and the two projections of $\Gamma_q$ give unconditional positive matricial coupling.

The construction is diagrammatic.  Transported Shilov ideals have a minimal join completion
$\mathcal M$ on every oriented semisimplicial quotient diagram.  With $\mathcal A$ the universal
cover coefficient system and $\mathcal Q=\mathcal A/\mathcal M$, there is a natural long exact
sequence
\[
 \cdots\longrightarrow H^k(\mathcal A)\longrightarrow H^k(\mathcal Q)
 \longrightarrow H^{k+1}(\mathcal M)\longrightarrow H^{k+1}(\mathcal A)
 \longrightarrow\cdots ,
\]
and $H^0(\mathcal Q)$ is exactly the algebra of compatible envelope correlations.  On an inverse
path this becomes
\[
 0\longrightarrow(\varprojlim A_n)/J_{\rm loc}\longrightarrow R_\Gamma
 \longrightarrow H^1(\mathcal M)\longrightarrow0.
\]
The usual $\varprojlim^1$ law and countable continuity of local $C^*$-envelopes are the
zero-escape specialization.

The obstruction is unrestricted.  Every unital $C^*$-algebra $D$ occurs in a split two-edge path
simultaneously as a deficit, an escape algebra, and the target-vertical ambiguity created by
relational composition; the migration group $H^1(\mathcal M)$ is canonically its underlying
involutive Banach space.  The composite quotient nevertheless induces an envelope epimorphism.
Hence the cohomology is not a reformulation of one-arrow nonfunctoriality, while its
$C^*$-algebraic edge faces retain arbitrary noncommutative structure.
\end{abstract}

\maketitle

\section{Introduction}\label{sec:intro}

The $C^*$-envelope of an operator system is its minimal generated $C^*$-cover
\cite{Hamana1979,Paulsen2002}.  Boundary representations recover its Shilov ideal
\cite{Arveson1969,Arveson2008,DritschelMcCullough2005}.  Minimality makes the envelope
well behaved under complete order embeddings, but not under quotients.  A complete order
quotient has positive lifts at every matrix level \cite{Han2011,KPTT2013}; nevertheless it need
not induce any map of $C^*$-envelopes.

For a quotient $q:S\twoheadrightarrow T$, universality gives a surjective $*$-homomorphism
\[
 \alpha:\Cu(S)\twoheadrightarrow\Cu(T).
\]
If $J_S,J_T$ are the Shilov ideals, an envelope map exists exactly when
\(\alpha(J_S)\subseteq J_T\).  The two failures of equality are
\[
 \Esc(q)=\frac{\alpha(J_S)}{\alpha(J_S)\cap J_T},
 \qquad
 \Def(q)=\frac{J_T}{\alpha(J_S)\cap J_T}.
\]
The first is boundary mass which escapes the target Shilov ideal; the second is target boundary
mass not reached from the source.  These algebras correctly diagnose one arrow, but they do not
by themselves supply an object that composes.

The central observation is that the sum, rather than either difference, is universal:
\[
 M_q=J_T+\alpha(J_S).
\]
It is the least target ideal forced simultaneously by the two endpoint boundaries.  Quotienting
by this join gives a common quotient $Q_q$ of the two envelopes, and their pullback over $Q_q$ is
the canonical envelope correlation $\Gamma_q$.  This produces three levels of structure.

\begin{enumerate}[label=\textup{(\roman*)}]
\item Every arrow becomes a closed subdirect $C^*$-relation.  Escape and deficit are the two
      quotient faces of its edge join.
\item Correlations compose.  The canonical assignment is normal oplax; inserting a vertex adds
      its transported Shilov ideal to the relational ambiguity.
\item On a diagram, the edge joins are the degree-one part of a minimal simplicial coefficient
      system.  Its long exact cohomology sequence measures all coherent failures of envelope
      descent.
\end{enumerate}

The third statement is the universal law.  It is formally adjacent to cohomology of diagrams,
but the coefficient system is forced by operator boundaries: the vertex ideals do not form a
coefficient system because escape is precisely the failure of transport.  The join completion is
the least enlargement that repairs every face map.  Ordinary finite ideal-lattice resolutions
\cite{Maszczyk2009} and Mayer--Vietoris constructions for $C^*$-ideals
\cite{Naarmann2019} concern several ideals in a fixed algebra.  Here the ideals live in changing
universal covers and are mixed by quotient transport.

There are other notions of noncommutative relation.  Weaver's quantum relations are weak-star
closed commutant bimodules \cite{Weaver2012}, and the recent two-algebra theory of Daws keeps that
operator-bimodule meaning \cite{Daws2026}.  Our correlations are subdirect $C^*$-products, hence
Goursat pullbacks; they retain quotient and ideal structure rather than support of a completely
positive map.  Restriction categories axiomatize partial maps \cite{CockettLack2002}, but a partial
map records only a domain of definition.  The join $M_q$ records the multivalued part as well, and
is therefore indispensable.

Section~\ref{sec:edge} proves the edge pullback and vertical exact sequence.
Section~\ref{sec:corr} constructs the correlation category and its oplax path law.
Section~\ref{sec:cohomology} proves the minimal join-completion theorem and long exact sequence.
Section~\ref{sec:path} identifies the path law and recovers inverse-limit continuity.
Section~\ref{sec:realization} proves that every unital $C^*$-algebra is realized by a two-edge
transmutation.  Section~\ref{sec:detection} gives the matricial and representation-theoretic
interpretations.

\section{Universal covers and boundary migration}\label{sec:prelim}

All operator systems are unital.  A map is a \emph{complete order quotient} if it induces a
complete order isomorphism from the Archimedean operator-system quotient by its kernel onto its
range.  In particular it is surjective and ucp.  A surjective ucp map with a ucp right inverse is
a complete order quotient.

For an operator system $S$, write
\[
 u_S:S\longrightarrow\Cu(S),\qquad
 \epsilon_S:\Cu(S)\longrightarrow\Ce(S)
\]
for the universal embedding and envelope quotient.  Its kernel $J_S$ is the Shilov ideal.
Every ucp map from $S$ to a $C^*$-algebra extends uniquely to a $*$-homomorphism from
$\Cu(S)$ after the target representation is fixed by the universal property.  Consequently a
complete order quotient $q:S\twoheadrightarrow T$ induces a unique unital $*$-homomorphism
\[
 \alpha_q:\Cu(S)\longrightarrow\Cu(T),\qquad
 \alpha_q u_S=u_Tq.
\]
It is onto because its range contains $u_T(T)$, which generates $\Cu(T)$.

\begin{definition}\label{def:exchange}
For a complete order quotient $q:S\twoheadrightarrow T$, define
\[
 I_q=\alpha_q(J_S)\cap J_T,\qquad
 \Esc(q)=\alpha_q(J_S)/I_q,\qquad
 \Def(q)=J_T/I_q.
\]
We call these the \emph{escape} and \emph{deficit} algebras.
\end{definition}

The elementary descent criterion will be used repeatedly.

\begin{proposition}[Envelope descent]\label{prop:descent}
The quotient $q$ extends to a surjective $*$-homomorphism
\(\bar q:\Ce(S)\twoheadrightarrow\Ce(T)\) if and only if $\Esc(q)=0$.
When it exists, $\bar q$ is unique.
\end{proposition}

\begin{proof}
The formula
\[
 \bar q(a+J_S)=\alpha_q(a)+J_T
\]
is well-defined exactly when $\alpha_q(J_S)\subseteq J_T$, equivalently when
$\Esc(q)=0$.  It is onto.  Conversely, two $*$-homomorphisms from $\Cu(S)$ into $\Ce(T)$
which agree on $u_S(S)$ agree everywhere.  Hence any extension satisfies
\(\bar q\epsilon_S=\epsilon_T\alpha_q\), forcing
\(\alpha_q(J_S)\subseteq J_T\).  Uniqueness follows from generation by $S$.
\end{proof}

\section{The edge pullback law}\label{sec:edge}

\begin{definition}\label{def:join-correlation}
For $q:S\twoheadrightarrow T$, put
\[
 M_q=J_T+\alpha_q(J_S)\triangleleft\Cu(T),\qquad
 Q_q=\Cu(T)/M_q,
\]
and define the \emph{envelope correlation}
\[
 \Gamma_q=
 \{(a+J_S,\alpha_q(a)+J_T):a\in\Cu(S)\}
 \subseteq\Ce(S)\oplus\Ce(T).
\]
\end{definition}

Finite sums of closed ideals in a $C^*$-algebra are closed, so $M_q$ is a closed ideal.

\begin{theorem}[Shilov edge pullback]\label{thm:pullback}
There are canonical surjective $*$-homomorphisms
\[
 \rho_S:\Ce(S)\twoheadrightarrow Q_q,\quad
 \rho_S(a+J_S)=\alpha_q(a)+M_q,
\]
\[
 \rho_T:\Ce(T)\twoheadrightarrow Q_q,\quad
 \rho_T(b+J_T)=b+M_q,
\]
and
\[
 \boxed{\Gamma_q=\Ce(S)\times_{Q_q}\Ce(T).}
\]
In particular $\Gamma_q$ is a closed subdirect $C^*$-product.

The ideal $M_q$ is the least ideal $N\triangleleft\Cu(T)$ containing $J_T$ for which the
transported map $\Ce(S)\to\Cu(T)/N$ is well-defined.  Its endpoint faces are
\[
 M_q/J_T\cong\Esc(q),\qquad
 M_q/\alpha_q(J_S)\cong\Def(q).
\]
\end{theorem}

\begin{proof}
Both endpoint maps are well-defined and onto.  A pair
\((a+J_S,b+J_T)\) belongs to their pullback exactly when
\[
 b-\alpha_q(a)\in J_T+\alpha_q(J_S).
\]
Write the difference as $j_T+\alpha_q(j_S)$.  The pair then equals
\[
 (a+j_S+J_S,\alpha_q(a+j_S)+J_T),
\]
which lies in $\Gamma_q$.  The reverse inclusion is immediate.

If $N$ permits both endpoint maps, then $J_T\subseteq N$ and well-definedness of the source map
forces $\alpha_q(J_S)\subseteq N$.  Thus $M_q\subseteq N$.  The two face formulas are the
second isomorphism theorem.
\end{proof}

\begin{proposition}[Least correlation property]\label{prop:least}
Inside \(\Ce(S)\oplus\Ce(T)\),
\[
 \Gamma_q=C^*\{(\epsilon_Su_S(s),\epsilon_Tu_T(qs)):s\in S\}.
\]
It is the least closed \(C^*\)-relation containing the operator-system graph of \(q\).
Equivalently, \(Q_q\) is the maximal common \(C^*\)-quotient of the endpoint envelopes on which
the two copies of that graph agree.
\end{proposition}

\begin{proof}
The defining image of \(\Cu(S)\) is generated by the image of
\(u_S(S)\), giving the displayed identity.  Any closed \(C^*\)-subalgebra containing the graph
therefore contains \(\Gamma_q\).  Its coordinate images contain the generating copies of \(S\)
and \(T\), so it is automatically subdirect.  The maximal-common-quotient assertion is the
minimality of \(M_q\) in Theorem~\ref{thm:pullback}, with the order reversed by quotienting.
\end{proof}

The pullback also identifies both vertical directions.

\begin{theorem}[Vertical extension]\label{thm:vertical}
Let $p_S,p_T$ be the coordinate maps of $\Gamma_q$, let
\[
 D_q=J_S\cap\alpha_q^{-1}(J_T),
\]
and let $\mathcal K_q$ be the image of $\ker\alpha_q$ in $\ker p_T$.  Then
\[
 \Gamma_q\cong\Cu(S)/D_q,\qquad
 \ker p_S\cong\Esc(q),
\]
and there is a natural short exact sequence
\[
 0\longrightarrow\mathcal K_q\longrightarrow\ker p_T
 \longrightarrow\Def(q)\longrightarrow0.
\]
\end{theorem}

\begin{proof}
The homomorphism
\[
 \theta_q:\Cu(S)\longrightarrow\Ce(S)\oplus\Ce(T),\qquad
 \theta_q(a)=(a+J_S,\alpha_q(a)+J_T)
\]
has kernel $D_q$.  Its image is closed, being a $C^*$-algebra quotient, and equals
$\Gamma_q$.  Its source-vertical kernel is
\[
 \{(0,\alpha_q(j)+J_T):j\in J_S\}
 \cong\alpha_q(J_S)/I_q=\Esc(q).
\]
The opposite vertical kernel is $\alpha_q^{-1}(J_T)/D_q$.  On it define
\[
 [a]\longmapsto\alpha_q(a)+I_q\in\Def(q).
\]
The map is well-defined and onto.  If $\alpha_q(a)\in I_q$, choose $j\in J_S$ with
$\alpha_q(a)=\alpha_q(j)$.  Then $j\in D_q$ and $[a]=[a-j]$, with
$a-j\in\ker\alpha_q$.  This identifies the kernel with $\mathcal K_q$.
\end{proof}

\begin{corollary}\label{cor:graph}
The correlation $\Gamma_q$ is the graph of a $*$-homomorphism
\(\Ce(S)\to\Ce(T)\) if and only if $\Esc(q)=0$.  Hence correlations extend the envelope
assignment to every complete order quotient, and the graph locus is its maximal functorial
subcategory.
\end{corollary}

\begin{proof}
A subdirect relation is the graph of a map from its first coordinate exactly when its first
coordinate projection is injective.  Apply Theorem~\ref{thm:vertical} and
Proposition~\ref{prop:descent}.
\end{proof}

\section{Correlation composition and transported joins}\label{sec:corr}

\begin{definition}
A $C^*$-\emph{correlation} from $A$ to $B$ is a closed $C^*$-subalgebra
$R\subseteq A\oplus B$ whose two coordinate maps are onto.  For correlations
$R\subseteq A\oplus B$ and $T\subseteq B\oplus C$, define
\[
 T\circ R=\{(a,c):\text{there is }b\in B\text{ with }(a,b)\in R,\ (b,c)\in T\}.
\]
\end{definition}

\begin{proposition}\label{prop:corr-category}
Correlations, with the preceding composition and diagonal identities, form a category $\Corr$.
Inclusions of correlations make it locally ordered, and composition is monotone in both
variables.
\end{proposition}

\begin{proof}
The fiber product
\[
 R\times_B T=\{(r,t)\in R\oplus T:p_B(r)=p_B(t)\}
\]
is a $C^*$-algebra.  Its image under the endpoint $*$-homomorphism into $A\oplus C$ is closed
and is $T\circ R$.  Subdirectness follows by lifting successively through the two coordinate
epimorphisms.  Associativity is ordinary associativity of relational composition.  The remaining
claims are immediate.
\end{proof}

Let $\OSq$ denote the category of operator systems and complete order quotients, with arrows
oriented from source to target.

\begin{theorem}[Oplax envelope law]\label{thm:oplax}
The assignments
\[
 S\longmapsto\Ce(S),\qquad q\longmapsto\Gamma_q
\]
form a normal oplax functor from $\OSq$ to the locally ordered category $\Corr$:
\[
 \Gamma_{rq}\subseteq\Gamma_r\circ\Gamma_q,\qquad
 \Gamma_{\id_S}=\{(x,x):x\in\Ce(S)\}.
\]
\end{theorem}

\begin{proof}
Write $\alpha,\beta$ for the universal epimorphisms of $q,r$.  Every element of
$\Gamma_{rq}$ has the form
\[
 (a+J_S,\beta\alpha(a)+J_U).
\]
Its middle coordinate $\alpha(a)+J_T$ witnesses membership in
$\Gamma_r\circ\Gamma_q$.  The identity assertion is immediate.
\end{proof}

The inclusion may be strict.  Its exact size along a path is an ideal join.

\begin{theorem}[Finite path law]\label{thm:pathjoin}
Consider a finite path
\[
 S_0\xrightarrow{q_1}S_1\xrightarrow{q_2}\cdots\xrightarrow{q_n}S_n,
\]
write $\alpha_{j:i}:A_i\to A_j$ for the composite universal epimorphism, and put
$J_i=J_{S_i}$.  If $P\subseteq\{0,\ldots,n\}$ is an ordered set of visited vertices containing
$0,n$, then the composed correlation through $P$ is
\[
 \left\{\left(a+J_0,\
 \alpha_{n:0}(a)+
 \sum_{i\in P\setminus\{0,n\}}\alpha_{n:i}(j_i)+J_n\right):
 a\in A_0,\ j_i\in J_i\right\}.
\]
Its target-vertical ideal is
\[
 \boxed{\operatorname{Amb}(P)=
 \frac{J_n+\sum_{i\in P\setminus\{n\}}\alpha_{n:i}(J_i)}{J_n}.}
\]
Consequently refinement of a path acts by ideal join.
\end{theorem}

\begin{proof}
For two edges, equality of the middle envelope coordinates says that the two middle
representatives differ by an element of $J_1$.  Transporting that element gives the displayed
formula.  Induction adds one transported intermediate ideal at each step.  For a fixed source
class, the representative $a$ may also vary through $J_0$, giving the vertical ideal.
\end{proof}

\begin{corollary}[Right-exact deficit composition]\label{cor:defcomp}
Let \(S\xrightarrow{q}T\xrightarrow{r}U\) be zero-escape complete order quotients, with universal
maps \(\alpha,\beta\).  Then
\[
 \Def(q)\xrightarrow{\beta_*}\Def(rq)\longrightarrow\Def(r)\longrightarrow0
\]
is exact, where
\[
 \ker\beta_*=
 \frac{\alpha(J_S)+(J_T\cap\ker\beta)}{\alpha(J_S)}.
\]
Thus it becomes short exact on the left after quotienting by the displayed kernel.
\end{corollary}

\begin{proof}
Zero escape gives the filtration
\[
 \beta\alpha(J_S)\subseteq\beta(J_T)\subseteq J_U.
\]
The three deficits are
\[
 J_T/\alpha(J_S),\qquad J_U/\beta\alpha(J_S),\qquad J_U/\beta(J_T).
\]
The quotient sequence of the filtration is exact in the middle and on the right.  The kernel of
the first induced map consists of classes represented by
\(J_T\cap\beta^{-1}(\beta\alpha(J_S))=
\alpha(J_S)+(J_T\cap\ker\beta)\).
\end{proof}

\begin{remark}
The assignment is generally oplax, not pseudo-functorial.  Strictness is not a defect of the
correlation category: Theorem~\ref{thm:realization} realizes arbitrary noncommutative strictness
while the direct composite is again a graph.
\end{remark}

\section{The simplicial Shilov coefficient system}\label{sec:cohomology}

We use algebraic cohomology of coefficient systems.  This avoids an unnecessary closed-range
assumption; topological refinements are discussed in Remark~\ref{rem:topology}.

\begin{definition}
An \emph{oriented semisimplicial quotient diagram} consists of an oriented semisimplicial set $K$,
an operator system $S_i$ at every vertex, and compatible complete order quotients along the
ordered vertex maps of every simplex.  Thus an ordered simplex
\[
 \sigma=[i_0,\ldots,i_p]
\]
has compatible universal epimorphisms
\(\alpha_{l:k}:A_{i_k}\twoheadrightarrow A_{i_l}\) for $k\leq l$, with
\(\alpha_{m:l}\alpha_{l:k}=\alpha_{m:k}\).
\end{definition}

The nerve of any quotient category is such a diagram.  A chain of specified arrows may instead
be regarded as a one-dimensional semisimplicial diagram; this distinction matters when one wants
the path complex rather than the full nerve.

For an ordered simplex $\sigma=[i_0,\ldots,i_p]$, define
\[
 \mathcal A(\sigma)=A_{i_p},\qquad
 \mathcal M(\sigma)=\sum_{k=0}^p\alpha_{p:k}(J_{i_k}),\qquad
 \mathcal Q(\sigma)=\mathcal A(\sigma)/\mathcal M(\sigma).
\]
For a nonterminal face, the terminal algebra is unchanged and the face join includes into the
simplex join.  For the terminal face, use the last universal epimorphism.

\begin{theorem}[Minimal Shilov join completion]\label{thm:joincompletion}
The preceding face maps make $\mathcal M\subseteq\mathcal A$ and $\mathcal Q$ coefficient
systems on $K$.  They fit into a short exact sequence of cochain complexes
\[
 0\longrightarrow C^\bullet(K;\mathcal M)
 \longrightarrow C^\bullet(K;\mathcal A)
 \longrightarrow C^\bullet(K;\mathcal Q)\longrightarrow0.
\]
The coefficient system $\mathcal M$ is minimal: if
$\mathcal N\subseteq\mathcal A$ is any coefficient subsystem with
$J_i\subseteq\mathcal N(i)$ at every vertex, then
\[
 \mathcal M(\sigma)\subseteq\mathcal N(\sigma)
\]
for every simplex.  Equivalently, $\mathcal Q$ is the largest quotient coefficient system whose
vertex values are the $C^*$-envelopes.
\end{theorem}

\begin{proof}
If the omitted vertex of a face is not terminal, the face join omits one summand of the simplex
join.  If it is terminal, the last epimorphism sends
\[
 \sum_{k=0}^{p-1}\alpha_{p-1:k}(J_{i_k})
 \quad\text{into}\quad
 \sum_{k=0}^{p}\alpha_{p:k}(J_{i_k}).
\]
Thus all coface maps are well-defined.  The coface identities away from the terminal face are
identities of ideal inclusions; those involving terminal faces follow from
\(\alpha_{m:l}\alpha_{l:k}=\alpha_{m:k}\).  Hence the alternating coface differential squares
to zero.

For each simplex,
\[
 0\to\mathcal M(\sigma)\to\mathcal A(\sigma)\to\mathcal Q(\sigma)\to0
\]
is exact.  Products of surjective linear maps are surjective, so taking products over all
$p$-simplices gives a short exact sequence in every degree, compatible with the differentials.

For minimality, include the vertex face $i_k$ into $\sigma$.  Coefficient compatibility forces
the transported image $\alpha_{p:k}(J_{i_k})$ into $\mathcal N(\sigma)$.  Summing over $k$
gives the result.
\end{proof}

\begin{theorem}[Shilov correlation long exact sequence]\label{thm:les}
Every oriented semisimplicial quotient diagram has a natural long exact sequence
\[
 \cdots\to H^k(K;\mathcal M)\to H^k(K;\mathcal A)
 \to H^k(K;\mathcal Q)\xrightarrow{\partial}
 H^{k+1}(K;\mathcal M)\to H^{k+1}(K;\mathcal A)\to\cdots.
\]
Moreover,
\[
 H^0(K;\mathcal Q)=
 \{(x_i)\in\prod_i\Ce(S_i):(x_i,x_j)\in\Gamma_f
 \text{ for every edge }f:i\to j\}.
\]
\end{theorem}

\begin{proof}
The long exact sequence is the standard connecting sequence of
Theorem~\ref{thm:joincompletion}.  On an edge $f:i\to j$,
\[
 \mathcal Q(f)=A_j/(J_j+\alpha_f(J_i))=Q_f.
\]
The degree-zero cocycle equation says that the two vertex classes have equal images in $Q_f$.
By Theorem~\ref{thm:pullback}, this is equivalent to membership in $\Gamma_f$.
\end{proof}

\begin{definition}\label{def:migration}
The boundary-specific migration group of a semisimplicial quotient diagram is
\[
 \operatorname{Mig}^{k+1}(K)=\operatorname{im}\!\left[
 H^k(K;\mathcal Q)\xrightarrow{\partial}H^{k+1}(K;\mathcal M)\right].
\]
By exactness,
\[
 \operatorname{Mig}^{k+1}(K)=
 \ker\!\left[H^{k+1}(K;\mathcal M)\longrightarrow
 H^{k+1}(K;\mathcal A)\right].
\]
\end{definition}

\begin{remark}
The distinction in Definition~\ref{def:migration} is essential in dimensions above one.
\(H^{k+1}(K;\mathcal M)\) may contain ordinary cohomology of the indexing complex even for a
constant identity diagram.  Only \(\operatorname{Mig}^{k+1}(K)\) measures failure of coherent
envelope descent.  On a one-dimensional inverse path,
\(H^1(K;\mathcal A)=0\), so \(\operatorname{Mig}^1(K)=H^1(K;\mathcal M)\).
\end{remark}

\begin{remark}\label{rem:topology}
Theorem~\ref{thm:les} is unconditional in complex vector spaces and in involutive vector spaces.
For infinite diagrams, quotient cohomology need not be Hausdorff unless the coboundary range is
closed.  No topological exactness is claimed without that hypothesis.  The principal
applications below are Hausdorff: the zero-escape local-envelope group vanishes, while the
realization theorem identifies its group with a prescribed Banach space.
\end{remark}

\section{Paths, local envelopes, and derived acyclicity}\label{sec:path}

Consider an inverse path
\[
 \cdots\longrightarrow S_2\xrightarrow{q_1}S_1\xrightarrow{q_0}S_0,
\]
put $A_n=\Cu(S_n)$, and write
\(\alpha_n:A_{n+1}\twoheadrightarrow A_n\).  Its one-dimensional coefficient complex is
\[
 C^0(\mathcal M)=\prod_nJ_n,\qquad
 C^1(\mathcal M)=\prod_nM_n,\qquad
 M_n=J_n+\alpha_n(J_{n+1}),
\]
\[
 d_{\mathcal M}((j_n))=(j_n-\alpha_n(j_{n+1}))_n.
\]
Define
\[
 R_\Gamma=
 \{(x_n)\in\prod_n\Ce(S_n):(x_{n+1},x_n)\in\Gamma_{q_n}\},
\]
\[
 A=\varprojlim A_n,\qquad
 J_{\rm loc}=\{(a_n)\in A:a_n\in J_n\text{ for every }n\}.
\]

\begin{theorem}[Migration exact sequence]\label{thm:migration}
There is a natural exact sequence
\[
 \boxed{
 0\longrightarrow A/J_{\rm loc}\longrightarrow R_\Gamma
 \xrightarrow{\delta}H^1(\mathcal M)\longrightarrow0.}
\]
For $x=(x_n)\in R_\Gamma$, choose representatives $a_n\in A_n$.  Then
\[
 \delta(x)=
 [(a_n-\alpha_n(a_{n+1}))_n].
\]
\end{theorem}

\begin{proof}
The relation condition and Theorem~\ref{thm:pullback} imply
\[
 a_n-\alpha_n(a_{n+1})\in M_n.
\]
Changing $a_n$ by $j_n\in J_n$ changes the mismatch by
$d_{\mathcal M}((j_n))$, so $\delta$ is well-defined.

If the mismatch class is zero, then
\[
 a_n-\alpha_n(a_{n+1})=j_n-\alpha_n(j_{n+1})
\]
for some $j_n\in J_n$.  The corrected representatives $a_n-j_n$ are coherent and represent the
same envelope classes.  This proves exactness at $R_\Gamma$; exactness on the left is immediate.

Given $c=(c_n)\in\prod M_n$, take $a_0=0$ and choose recursively
\[
 \alpha_n(a_{n+1})=a_n-c_n.
\]
Every choice is possible because $\alpha_n$ is onto.  The classes $a_n+J_n$ form a relational
family with mismatch $c$, so $\delta$ is onto.
\end{proof}

The edge coefficient simultaneously contains the exchange pair:
\[
 M_n/J_n\cong\Esc(q_n),\qquad
 M_n/\alpha_n(J_{n+1})\cong\Def(q_n).
\]
Thus the one-arrow invariants are quotient stalks, while $H^1(\mathcal M)$ is their coherent
global obstruction.

\begin{corollary}[Zero-escape specialization]\label{cor:roos}
If every $q_n$ has zero escape, then $M_n=J_n$, every $\Gamma_{q_n}$ is the graph of the envelope
epimorphism, and Theorem~\ref{thm:migration} becomes
\[
 0\longrightarrow A/J_{\rm loc}\longrightarrow
 \varprojlim(A_n/J_n)\longrightarrow
 \varprojlim\nolimits^1J_n\longrightarrow0,
\]
where
\[
 \varprojlim\nolimits^1J_n=
 \coker\!\left[
 \prod J_n\longrightarrow\prod J_n,\quad
 (j_n)\mapsto(j_n-\alpha_n(j_{n+1}))
 \right].
\]
\end{corollary}

\begin{corollary}[Countable envelope continuity]\label{cor:continuity}
Assume the inverse path is countable and defines a Fr\'echet local operator system
$S=\varprojlim S_n$.  If every link has zero escape, then
\[
 C_{e,\mathrm{loc}}^*(S)\cong\varprojlim\Ce(S_n),
\qquad
 \varprojlim\nolimits^1J_n=0.
\]
\end{corollary}

\begin{proof}
First identify the local Shilov ideal.  The set \(J_{\rm loc}\) is a closed ideal of \(A\).
For \(s\in M_r(S)\), its \(n\)th quotient seminorm modulo \(J_{\rm loc}\) is bounded above by
\(\|s_n\|\) and below by the quotient norm modulo \(M_r(J_n)\).  Since \(J_n\) is boundary for
\(S_n\), both bounds are \(\|s_n\|\).  Thus \(J_{\rm loc}\) is a local boundary ideal.

Conversely, let \(H\triangleleft A\) be a local boundary ideal.  The closed ideal
\(\overline{p_n(H)}\triangleleft A_n\) is boundary for \(S_n\): every matrix over \(S_n\) lifts
to a matrix over \(S\) by recursive use of the surjective bonding maps, and the \(n\)th local
quotient seminorm is the quotient norm by \(\overline{p_n(H)}\).  Maximality of \(J_n\) gives
\(\overline{p_n(H)}\subseteq J_n\) for every \(n\), hence \(H\subseteq J_{\rm loc}\).  Therefore
\(A/J_{\rm loc}=C_{e,\mathrm{loc}}^*(S)\) by local Shilov theory
\cite[Theorems~3.16 and~3.18]{Joita2025}.

Now put \(B=\varprojlim\Ce(S_n)\).  Each projection \(S\to S_n\) is onto: lift recursively above
stage \(n\) and define lower coordinates by composition.  Since \(S_n\) generates \(\Ce(S_n)\),
polynomial approximation at coordinate \(n\) shows that \(S\) locally generates \(B\).
If \(L\triangleleft B\) is a local boundary ideal, then
\(\overline{p_n(L)}\) is boundary for \(S_n\subseteq\Ce(S_n)\), whose Shilov ideal is zero.
Thus every coordinate image of \(L\) is zero and \(L=0\).

Theorem~\ref{thm:migration} gives the canonical injective local
\(*\)-homomorphism \(\kappa:A/J_{\rm loc}\to B\) fixing \(S\).  Both generated local covers have
zero local Shilov ideal, so local-envelope uniqueness
\cite[Theorem~3.18]{Joita2025} implements the identity on \(S\) by a local
\(*\)-isomorphism.  It must equal \(\kappa\) because \(S\) locally generates both sides.
Hence \(\kappa\) is onto, and Corollary~\ref{cor:roos} forces
\(\varprojlim^1J_n=0\).
\end{proof}

\begin{remark}
The vanishing in Corollary~\ref{cor:continuity} is not a pointwise Mittag--Leffler statement.
The maps $J_{n+1}\to J_n$ need not be onto; equivalently, link deficits may be nonzero.  Local
envelope minimality kills their coherent degree-one class rather than each edge quotient.
\end{remark}

\section{Universal two-edge transmutation}\label{sec:realization}

We now show that migration cohomology is already unrestricted on the smallest path that can
distinguish direct from staged transport.

\begin{theorem}[Universal transmutation]\label{thm:realization}
For every unital $C^*$-algebra $D$, there are split complete order quotients
\[
 S_2\xrightarrow{q_1}S_1\xrightarrow{q_0}S_0
\]
such that
\[
 \Def(q_1)\cong D,\qquad \Esc(q_0)\cong D,
\]
$q_1$ and $q_0q_1$ have zero escape, and
\[
 \Gamma_{q_0}\circ\Gamma_{q_1}\supsetneq\Gamma_{q_0q_1}.
\]
More precisely,
\[
 \frac{\Gamma_{q_0}\circ\Gamma_{q_1}}{\Gamma_{q_0q_1}}
 \cong H^1(\mathcal M)\cong D
\]
as involutive vector spaces, where the first quotient is by the closed graph subspace.  Thus one
and the same $D$ occurs as deficit, escape, strict oplax composition excess, and migration
cohomology, although the composite again induces an envelope epimorphism.
\end{theorem}

\begin{proof}
Put
\[
 S_0=D,\qquad S_1=D\oplus D,\qquad
 q_0(a,b)=\tfrac12(a+b).
\]
The diagonal map is a ucp section, so $q_0$ is a split complete order quotient.  Let
$v=(1,-1)$.  The correlation $\Gamma_{q_0}$ contains $(v,0)$ and therefore
$(1,0)=(v,0)^*(v,0)$.  Subtracting from the unit gives $(0,1)$.  Multiplying the generating
graph by these two central projections yields both summands, and hence
\[
 \Gamma_{q_0}=(D\oplus D)\oplus D.
\]
Its target-vertical ideal is $D$, so Theorem~\ref{thm:vertical} gives
$\Esc(q_0)\cong D$.

Let
\[
 K=C([0,1],D),\qquad I=\{f\in K:f(0)=0\},\qquad
 \lambda(t)=\tfrac12+\tfrac14t,
\]
and define
\[
 \Phi:D\oplus D\longrightarrow K,\qquad
 \Phi(a,b)(t)=\lambda(t)a+(1-\lambda(t))b.
\]
The map is ucp: pointwise it is a convex combination of the coordinate maps.  Put
\[
 B=(D\oplus D)\oplus K,\qquad r(x)=(x,\Phi(x)),
\]
\[
 S_2=r(S_1)+(0\oplus I)\subseteq B,\qquad
 q_1(r(x)+(0,i))=x.
\]
The sum is closed: the first coordinate recovers $x$ contractively and then recovers
$i$ with a uniform bound.  First projection and $r$ are a ucp quotient and a ucp section.
Thus $q_1$ is a split complete order quotient with kernel $0\oplus I$.

We compute the envelope of $S_2$.  Since $\Phi(v)(t)=t/2$,
\[
 r(1)-r(v)^*r(v)=(0,1-t^2/4).
\]
Its spectrum in $B$ is contained in $\{0\}\cup[3/4,1]$.  Continuous functional calculus gives
the central projection $0\oplus1_K$ in $C^*(S_2)$.  The second coordinates $\Phi(S_1)$ contain
the constant functions from $D$ and the central function $t$, so they generate $K$.  It follows
that $C^*(S_2)=B$.

The graph subsystem $r(S_1)$ has Shilov ideal $0\oplus K$: quotient by this ideal is the minimal
cover $D\oplus D$.  Any boundary ideal for the larger system $S_2$ is therefore contained in
$0\oplus K$.  Complete isometry on the contained subspace $0\oplus I$ forces its intersection
with $0\oplus I$ to be zero.  The ideal $I$ is essential in $K$.  Indeed, if a nonzero ideal of
$K$ contains $h\ne0$, continuity gives $t>0$ with $h(t)\ne0$; multiplying by a scalar function
which vanishes at $0$ and equals $1$ at $t$ produces a nonzero element of its intersection with
$I$.  Hence the boundary ideal of $S_2$ is zero:
\[
 \Ce(S_2)=B.
\]

The envelope map of $q_1$ is first projection $p:B\to D\oplus D$, so $\Esc(q_1)=0$.
Its kernel is $0\oplus K$, while the ideal generated by the operator-system kernel is
$0\oplus I$.  The standard kernel form of deficit gives
\[
 \Def(q_1)\cong(0\oplus K)/(0\oplus I)\cong D.
\]
For completeness, that kernel formula follows from
\[
 \Def(q)\cong
 \ker\bar q\big/\Ideal_{\Ce(S)}(\ker q)
\]
whenever $\Esc(q)=0$: pull $\ker\bar q$ back to $\Cu(S)$ and use
\(\ker\alpha_q=\Ideal_{\Cu(S)}(u_S(\ker q))\).

The composite extends to the envelope epimorphism
\[
 \pi:B\longrightarrow D,\qquad
 \pi((a,b),f)=f(0),
\]
because $\Phi(a,b)(0)=\tfrac12(a+b)$ and $I(0)=0$.  Hence
$\Esc(q_0q_1)=0$ and $\Gamma_{q_0q_1}=\graph(\pi)$.  Since
$\Gamma_{q_1}=\graph(p)$ and $\Gamma_{q_0}$ is full,
\[
\Gamma_{q_0}\circ\Gamma_{q_1}=B\oplus D.
\]
The target-vertical ideal of this staged correlation is \(0\oplus D\), whereas the direct graph
has zero target-vertical ideal.  Thus the full \(C^*\)-algebraic ambiguity created by staging is
canonically \(D\).
The map
\[
 B\oplus D\longrightarrow D,\qquad (b,d)\longmapsto d-\pi(b)
\]
has kernel $\graph(\pi)$ and is onto, proving the composition-excess assertion.
Theorem~\ref{thm:migration}, applied to this finite path, identifies the same quotient with
$H^1(\mathcal M)$.
\end{proof}

\begin{corollary}\label{cor:no-bounds}
Split complete-order lifting imposes no dimensional, separability, exactness, nuclearity,
type~I, ideal-lattice, representation-theoretic, or $K$-theoretic restriction on the
$C^*$-algebraic escape, deficit, and composition-ambiguity faces.  The associated
$H^1(\mathcal M)$ retains the underlying involutive Banach space.
\end{corollary}

\begin{proof}
Choose $D$ with the desired property in Theorem~\ref{thm:realization}.
\end{proof}

\section{Matricial lifting and faithful boundary detection}\label{sec:detection}

The correlation does not merely retain ideal arithmetic.  It provides a two-sided positive
coupling at all matrix levels.

\begin{theorem}[Unconditional matricial coupling]\label{thm:matrix}
Let $q:S\twoheadrightarrow T$.  For every $r\geq1$, both maps
\[
 p_S^{(r)}:M_r(\Gamma_q)_+\twoheadrightarrow M_r(\Ce(S))_+,\qquad
 p_T^{(r)}:M_r(\Gamma_q)_+\twoheadrightarrow M_r(\Ce(T))_+
\]
are onto.  Positive lifts may be chosen with quotient-norm control.
\end{theorem}

\begin{proof}
The coordinate maps of $\Gamma_q$ are surjective $*$-homomorphisms by
Theorem~\ref{thm:pullback}.  Matrix amplification preserves this property, and positive elements
lift positively through $C^*$-quotients with the standard quotient-norm estimate.
\end{proof}

This should not be confused with lifting \emph{inside} Shilov ideals.

\begin{proposition}[Deficit and Shilov lifting]\label{prop:shilovlift}
Assume $\Esc(q)=0$.  Then the following are equivalent:
\begin{enumerate}[label=\textup{(\roman*)}]
\item $\Def(q)=0$;
\item $\alpha_q(J_S)=J_T$;
\item for every $r\geq1$, every $b\in M_r(J_T)_+$ and every $\varepsilon>0$, there is
      $a\in M_r(J_S)_+$ with
      \[
      \alpha_q^{(r)}(a)=b,\qquad \|a\|\leq\|b\|+\varepsilon.
      \]
\end{enumerate}
\end{proposition}

\begin{proof}
Under zero escape, $\alpha_q(J_S)\subseteq J_T$, so
\(\Def(q)=J_T/\alpha_q(J_S)\).  This proves the first equivalence.  A surjective
$*$-homomorphism of $C^*$-algebras has positive matrix lifts with quotient-norm control.
Conversely, scalar positive lifting implies surjectivity because every element of a $C^*$-algebra
is a linear combination of positive elements.
\end{proof}

For separable operator systems, direct sums of boundary representations detect the edge join.

\begin{proposition}[Faithful boundary detection]\label{prop:detection}
Assume $T$ is separable and let $\rho_T$ be a direct sum of sufficiently many boundary
representations of $T$ so that
\(\ker\rho_T=J_T\) on $\Cu(T)$.  Then
\[
 \rho_T(\alpha_q(J_S))\cong\Esc(q).
\]
Thus escape is exactly the transported source boundary mass seen by a faithful positive
representation of the target envelope.  The unseen complementary face of the same edge join is
$\Def(q)=M_q/\alpha_q(J_S)$.
\end{proposition}

\begin{proof}
The kernel of $\rho_T$ on $\alpha_q(J_S)$ is
\(\alpha_q(J_S)\cap J_T\).  Apply the first isomorphism theorem.
\end{proof}

\section{Consequences and scope}\label{sec:scope}

The algebraic path law has been checked independently in Lean~4 against the repository's fixed
mathlib toolchain.  The module
\[
 \texttt{GaloisForLFunctions.HTGeometry.ShilovEnvelopeCorrelationLaw}
\]
defines edge joins, relational representative classes, the migration difference and connecting
map, and proves injectivity, exactness at the relational term, surjectivity under ambient
surjective bonds, and the combined short exact theorem.  The proofs introduce no project axiom,
placeholder, opaque declaration, or unsafe definition.  Their only printed logical dependencies
are \texttt{propext}, \texttt{Classical.choice}, and \texttt{Quot.sound}.

This formal statement is deliberately algebraic.  It does not replace complete positivity,
Shilov maximality, functional calculus, essentiality, or local \(C^*\)-envelope uniqueness by
surrogate assumptions.  Those analytic arguments remain the paper proofs above.

The preceding results organize the standard phenomena as follows.

\begin{enumerate}[label=\textup{(\arabic*)}]
\item Failure of an envelope map is the multivalued vertical ideal of $\Gamma_q$
      (Theorem~\ref{thm:vertical}).
\item Kernel nonexactness is the opposite vertical quotient
      (Theorem~\ref{thm:vertical} and the kernel formula in
      Theorem~\ref{thm:realization}).
\item Positive matricial lifting is supplied unconditionally by the two correlation projections;
      Shilov-ideal lifting is the zero-deficit specialization
      (Theorem~\ref{thm:matrix} and Proposition~\ref{prop:shilovlift}).
\item Composition is ideal join, and the earlier right-exact deficit sequence is the quotient
      sequence of the three-term filtration
      \[
      \beta\alpha(J_S)\subseteq\beta(J_T)\subseteq J_U
      \]
      on the zero-escape locus.
\item The Roos sequence, local envelope continuity, and derived Shilov acyclicity are the
      one-dimensional zero-escape specialization of Theorem~\ref{thm:les}.
\item Boundary representations detect the vertex ideals, hence their minimal edge and simplex
      joins (Proposition~\ref{prop:detection}).
\end{enumerate}

The law does not assert that $H^k(\mathcal M)$ always vanishes.  Theorem~\ref{thm:realization}
shows the opposite in degree one.  Nor does it identify algebraic cohomology with a $C^*$-algebra
in general: a coboundary range need not be an ideal or be closed.  What is canonical is the
coefficient system, its exact sequence, and the compatible-correlation algebra in degree zero.

The universal statement can be summarized in one sentence:
\begin{quote}
\emph{The Shilov boundary of a quotient diagram is the minimal join coefficient system whose
quotient turns every $C^*$-envelope arrow into a correlation; its cohomology is exactly the
obstruction to replacing those correlations by coherent envelope maps.}
\end{quote}

\sloppy
\begin{thebibliography}{99}

\bibitem{Arveson1969}
W.~Arveson,
\emph{Subalgebras of $C^*$-algebras},
Acta Math. \textbf{123} (1969), 141--224.

\bibitem{Arveson2008}
W.~Arveson,
\emph{The noncommutative Choquet boundary},
J. Amer. Math. Soc. \textbf{21} (2008), 1065--1084.

\bibitem{CockettLack2002}
J.~R.~B.~Cockett and S.~Lack,
\emph{Restriction categories I: categories of partial maps},
Theoret. Comput. Sci. \textbf{270} (2002), 223--259.

\bibitem{Daws2026}
M.~Daws,
\emph{Quantum relations in the general setting: composition and adjacency operators},
arXiv:2602.18361 (2026).

\bibitem{DritschelMcCullough2005}
M.~A.~Dritschel and S.~A.~McCullough,
\emph{Boundary representations for families of representations of operator algebras and spaces},
J. Operator Theory \textbf{53} (2005), 159--167.

\bibitem{Hamana1979}
M.~Hamana,
\emph{Injective envelopes of operator systems},
Publ. Res. Inst. Math. Sci. \textbf{15} (1979), 773--785.

\bibitem{Han2011}
K.~H.~Han,
\emph{On maximal tensor products and quotient maps of operator systems},
J. Math. Anal. Appl. \textbf{384} (2011), 375--386.

\bibitem{Joita2025}
M.~Joi\c{t}a,
\emph{The Shilov boundary for a local operator system},
Adv. Oper. Theory \textbf{10} (2025), no.~3,
\href{https://doi.org/10.1007/s43036-025-00459-5}{doi:10.1007/s43036-025-00459-5}.

\bibitem{KPTT2013}
A.~S.~Kavruk, V.~I.~Paulsen, I.~G.~Todorov, and M.~Tomforde,
\emph{Quotients, exactness, and nuclearity in the operator system category},
Adv. Math. \textbf{235} (2013), 321--360.

\bibitem{Maszczyk2009}
T.~Maszczyk,
\emph{Distributive lattices and cohomology},
arXiv:0811.3997v2 (2009).

\bibitem{Naarmann2019}
S.~Naarmann,
\emph{A Mayer--Vietoris spectral sequence for $C^*$-algebras and coarse geometry},
arXiv:1812.11442v2 (2019).

\bibitem{Paulsen2002}
V.~I.~Paulsen,
\emph{Completely bounded maps and operator algebras},
Cambridge Studies in Advanced Mathematics, vol.~78, Cambridge University Press, 2002.

\bibitem{Weaver2012}
N.~Weaver,
\emph{Quantum relations},
Mem. Amer. Math. Soc. \textbf{215} (2012), no.~1010.

\end{thebibliography}

\end{document}
